\section{What Becomes of Anubion?}


Clement takes advantage\footnote{\url{https://www.scribd.com/doc/13561289/-The-Authentic-Peter-The-Preaching-of-Simeon-Kefa-from-the-Journal-of-T-Flavius-Clemens-Clement-The-Clementine-Recognitions-Recognitions-of-Cl}} of the degeneration of Tradition to make an argument against Astrology:

\begin{quotex}
And this is why ignorant astrologers have invented to themselves the talk about climacterics as their refuge in uncertainties, as we showed fully yesterday.
\end{quotex}


This was a common Christian argument against ``the stars'' and the ``old ways'' in favor of Free Will. But what if man was ``more free'' than before, thanks to the Advent? In this case, there is really no conflict -- the laws of the universe have not altered, but the pattern has.

The source of Christian humility is rooted in the Archetypal nature of the Father:

\begin{quotex}
But no one ought to be ashamed of this, because there is no man who ought to profess that he knows all things; for there is only One who knows all things, even He who also made all things. For if our Master declared that He knew not the day and the hour whose signs even He foretold, and referred the whole to the Father, how shall we account it disgraceful to confess that we are ignorant of some things, since in this we have the example of our Master?
\end{quotex}


Christian charity is not a sacrifice, esoterically, but an enrichment of individual reward:

\begin{quotex}
Or do you not know that if all tribes, after they have heard from me the preaching of the truth, and have believed, would betake themselves to teaching, they would gain the greater kavod for me, if indeed you think me desirous of glory? For what so glorious as to prepare talmidim for Moshiach, not who shall be silent, and shall be saved alone, but who shall speak what they have learned, and shall do good to others?
\end{quotex}


The monotheistic objection to the cosmogonies of even the more ``intelligent'' Gentiles is arrived at, after dismissing Orpheus and Hesiod:

\begin{quotex}
that by such crimes and impieties they may teach men to imitate their dei whom they worship. For in these allegories what profit can there be to them? For although they are framed so as to be decent, yet no use is derived from them for worship, nor for amendment of morals... 
\end{quotex}


So the primary objection is a question of imitation, and one thinks of Carlyle's chapter on the French Revolution, of Saturn devouring his children:

\begin{quotex}
But since those things rather are true which the public baseness testifies, and concealments have been sought and fabricated by prudent men to excuse them by seemly speeches, therefore are they not only not prohibited, but even in the very mysteries figures are produced of Saturn devouring his sons, and of the boy hidden by the cymbals and drums of the Corybantes; and with respect to the mutilation of Saturn, what better proof of its truth could there be than that even his worshippers are mutilated, by a like miserable fate, in honor of their dei?
\end{quotex}


An interesting point is made by Peter, concerning General Law's allowing all pagan practices to degenerate in order to make men ready for full Truth:

\begin{quotex}
For the mind is with difficulty transferred from those things with which it has been imbued in early youth; and on this account, as I said, it has been affected by YHWH, that the substance of error should be both weak and base. But all other things also YHWH dispenses fitly and advantageously, although the method of the divine dispensation, as good, and the best possible, is not clear to us who are ignorant of the causes of things.
\end{quotex}


But this means that in the past, these very fables contained Truth, which was corrupted (over time and in time) by demons. It is not clear that Peter, the apostle, laid a clear emphasis at this point upon what is implicit in his words -- that mankind was once the receiver of natural revelation in paganism (a point he acknowledges elsewhere). The real question is, could Christianity itself (as a vehicle and religion) become subject to the same processes? The evidence we have today suggests that it will -- ``Decayed altars are inhabited by demons''.

Niceta explains the beauty of the allegories:

Then Niceta:

\begin{quotex}
The affair of the supper of the Elohim stands in this wise. They say that the banquet is the world, that the order of the Elohim sitting at table is the position of the heavenly bodies. Those whom Hesiod calls the first children of heaven and earth, of whom six were males and six females, they refer to the number of the twelve signs, which go round all the world. They say that the dishes of the banquet are the reasons and causes of things, sweet and desirable, which in the shape of inferences from the positions of the signs and the courses of the stars, explain how the world is ruled and governed. Yet they say these things exist after the free manner of a banquet, inasmuch as the mind of every one has the option whether he shall taste aught of this sort of knowledge, or whether he shall refrain; and as in a banquet no one is compelled, but every one is at liberty to eat, so also the manner of philosophizing depends upon the choice of the will. They say that discord is the lust of the flesh, which rises up against the purpose of the mind, and hinders the desire of philosophizing; and therefore they say that the time was that in which the marriage was celebrated..
\end{quotex}


Thus, at the core, paganism was compatible with the new Faith, in its esoteric mission, as a purification of the Temple of God (of which things, Jesus was a picture, in coming to the harlot Jerusalem and casting down the nation of the Judaizers).

Kefa (Simon Peter) confirms this:

Then Kefa, commending his statement, said:

\begin{quotex}
Ingenious men, as I perceive, take many counterfeits looking much like truth from the things which they read; and therefore great care is to be taken,that when the Mitzvot of Elohim is read, it be not read according to the understanding of our own mind.
\end{quotex}


Thus, the patristic emphasis on ``Tradition'', which was esotericism formalized as much as possible, but no more. It is even promised that the yoke will prove easy, the burden light.

\begin{quotex}
And concerning His prophetic foreknowledge and power, if any one,as I have said, wishes to receive clear proofs, let him come instantly and be alert to hear, and we shall give evident proofs by which he shall seem not only to hear the power of prophetic foreknowledge with his ears, but even to see it with his eyes and handle it with his hand; and when he has entertained a sure faith concerning Him, he will without any labor take upon him the yoke of zedekah and obedience; and so great sweetness will he perceive in it, that not only will he not find fault with any labor being in it, but will even desire something further to be added and imposed upon him.
\end{quotex}


The story ends with a plot twist, which I will not give away here -- suffice it to say that Shimon Magus shows back up, manages to transform someone's countenance, and has to be thwarted by the miracles, the wisdom, and even a pious lie of St. Peter. Interestingly enough, it appears that an astrologer named Anubion was an unbeliever, but a friend of Simon Peter.

We have reached the end of the Journal, although there is a letter from Peter to James the Just that we will finish with. What is there to learn? A great deal about the ``Genesis'' of the early Faith, the conditions of the ``Masses'', the secret doctrines of St. Peter (which are ``hid'' in the open, but require a good deal of apostolic training and explaining), and the possible rapprochement between sacred Traditions outside Christianity with a renewed zeal for the old ``Faith''.


\flrightit{Posted on 2012-09-07 by Logres}

